% ==========================================================================
%  Guia de Estudo — Prova 2 / Lista II
%  Macroeconomia I (PPGEco-UFES, 2026)
%  Prof. Ricardo Ramalhete Moreira
%  Rigidez Nominal e Modelos Novo-Keynesianos
% ==========================================================================
\documentclass[12pt, a4paper]{article}

% --------------------------------------------------------------------------
% Packages
% --------------------------------------------------------------------------
\usepackage[utf8]{inputenc}
\usepackage[T1]{fontenc}
\usepackage[brazil]{babel}
\usepackage{lmodern}
\usepackage{microtype}

\usepackage[top=2.5cm, bottom=2.5cm, left=3cm, right=3cm, headheight=15pt]{geometry}

\usepackage{amsmath, amssymb, amsthm}
\usepackage{mathtools}
\usepackage{bm}

\usepackage{graphicx}
\graphicspath{{lista_II/figures/}{figures/}}
\usepackage{float}
\usepackage[font=small, labelfont=bf, justification=justified]{caption}
\usepackage{subcaption}

\usepackage{booktabs}
\usepackage{array}
\usepackage{multicol}
\usepackage{xcolor}
\usepackage{hyperref}
\usepackage{enumitem}
\usepackage{tcolorbox}
\tcbuselibrary{breakable}

% --------------------------------------------------------------------------
% Cores
% --------------------------------------------------------------------------
\definecolor{cyanrbc}{HTML}{3b9dbd}
\definecolor{orangerbc}{HTML}{d97b39}
\definecolor{grayrbc}{HTML}{6b7280}
\definecolor{greenbox}{HTML}{166534}
\definecolor{greenbg}{HTML}{dcfce7}
\definecolor{boxbg}{HTML}{f5f0e8}
\definecolor{boxborder}{HTML}{3b9dbd}
\definecolor{provabg}{HTML}{fff7ed}
\definecolor{provaborder}{HTML}{d97b39}

% --------------------------------------------------------------------------
% Caixas
% --------------------------------------------------------------------------
\newtcolorbox{resultbox}[1][]{standard, colback=boxbg, colframe=boxborder,
  boxrule=1pt, arc=3pt, left=6pt, right=6pt, top=4pt, bottom=4pt,
  fonttitle=\bfseries\small, #1}
\newtcolorbox{provabox}[1][]{standard, colback=provabg, colframe=provaborder,
  boxrule=1.2pt, arc=3pt, left=6pt, right=6pt, top=4pt, bottom=4pt,
  fonttitle=\bfseries\small, #1}
\newtcolorbox{formulabox}[1][]{standard, colback=greenbg!30, colframe=greenbox!60,
  boxrule=0.8pt, arc=3pt, left=6pt, right=6pt, top=4pt, bottom=4pt,
  fonttitle=\bfseries\small, #1}

% --------------------------------------------------------------------------
% Comandos matemáticos
% --------------------------------------------------------------------------
\newcommand{\E}{\mathbb{E}}
\newcommand{\dif}{\mathrm{d}}
\newcommand{\pibar}{\bar{\pi}}
\newcommand{\xgap}{x_t}

% --------------------------------------------------------------------------
% Espaçamento e parágrafos
% --------------------------------------------------------------------------
\setlength{\parskip}{0.5\baselineskip}
\setlength{\parindent}{0pt}

% --------------------------------------------------------------------------
% Títulos
% --------------------------------------------------------------------------
\usepackage{titlesec}
\titleformat{\section}{\large\bfseries\color{cyanrbc}}{Questão \thesection.}{0.5em}{}[\vspace{-0.3em}\rule{\linewidth}{0.4pt}]
\titleformat{\subsection}{\normalsize\bfseries\color{grayrbc}}{}{0em}{}

% --------------------------------------------------------------------------
% Cabeçalho/rodapé
% --------------------------------------------------------------------------
\usepackage{fancyhdr}
\pagestyle{fancy}
\fancyhf{}
\lhead{\footnotesize\textcolor{grayrbc}{Guia P2 --- Macroeconomia I (PPGEco-UFES, 2026)}}
\rhead{\footnotesize\textcolor{grayrbc}{Rigidez Nominal \& NK --- Romer, cap.\ 6--7}}
\cfoot{\small\thepage}
\renewcommand{\headrulewidth}{0.4pt}

% ==========================================================================
\begin{document}
% ==========================================================================

% --------------------------------------------------------------------------
% Capa
% --------------------------------------------------------------------------
\begin{titlepage}
\begin{center}
\vspace*{1.5cm}
{\Large\bfseries Universidade Federal do Espírito Santo}\\[0.5em]
{\large Programa de Pós-Graduação em Economia (PPGEco-UFES)}\\[2em]

\rule{\linewidth}{1.2pt}\\[0.8em]
{\LARGE\bfseries\color{cyanrbc} Guia de Estudo --- Prova 2 / Lista II\\[0.4em]
\large Macroeconomia I --- 2026}\\[0.8em]
\rule{\linewidth}{1.2pt}

\vspace{1em}
{\large Prof.\ Ricardo Ramalhete Moreira}

\vspace{1.5em}
\begin{tcolorbox}[standard, colback=provabg, colframe=provaborder, width=0.88\textwidth,
  arc=4pt, boxrule=1.2pt, left=8pt, right=8pt, top=6pt, bottom=6pt]
\centering
\textbf{\textcolor{provaborder}{Atenção:}} Este guia foi elaborado com base na
\textbf{Lista II 2026} e cobre os modelos de
\textbf{Rigidez Nominal} (Romer, cap.\ 6) e \textbf{DSGE Novo-Keynesiano}
(Romer, cap.\ 7).
\end{tcolorbox}

\vspace{1.5em}
\begin{tabular}{ll}
\textbf{Questões 1--4:}  & Rigidez nominal, competição monopolística e preços \\[0.4em]
\textbf{Questões 5--7:}  & Informação incompleta e Crítica de Lucas \\[0.4em]
\textbf{Questões 8--10:} & Modelos Fischer, Taylor, Calvo e Christiano et al. \\[0.4em]
\textbf{Questão 11:}     & Modelo NK canônico --- IS, CPNK e Regra de Taylor
\end{tabular}

\vspace{1.5em}
\begin{tcolorbox}[standard, colback=boxbg, colframe=boxborder, width=0.88\textwidth,
  arc=4pt, boxrule=1pt, left=8pt, right=8pt, top=6pt, bottom=6pt]
\centering
\textbf{Referência principal:} \\[0.3em]
D.\ Romer, \textit{Advanced Macroeconomics}, 5.ª ed., McGraw-Hill (2019). \\
Capítulos 6 e 7.
\end{tcolorbox}

\vfill
{\small Soluções computadas com Python; código disponível em
\texttt{github.com/fcarva/macro}.}
\end{center}
\end{titlepage}

% --------------------------------------------------------------------------
% Formulário de Referência Rápida
% --------------------------------------------------------------------------
\newpage
\thispagestyle{fancy}

\begin{center}
{\large\bfseries\color{cyanrbc} Formulário de Referência Rápida --- P2}\\[0.3em]
\rule{\linewidth}{0.6pt}
\end{center}

\vspace{0.3em}
\begin{multicols}{2}

\textbf{\textcolor{cyanrbc}{BLOCO RIGIDEZ NOMINAL}}

\medskip
\textbf{Preço ótimo (competição monopolística):}
\[
P_i^* = \frac{\eta}{\eta-1}\,MC_i = \mu\,MC_i
\]

\textbf{Produto de equilíbrio (competição monop.):}
\[
Y^* = \left[\frac{\eta-1}{\eta}\right]^{1/(\gamma-1)}
\]

\textbf{Ganho de ajuste (menu costs):}
\[
G(p_{\text{dev}}) = \frac{\eta-1}{2}\,p_{\text{dev}}^2
\]

\textbf{Limiar de ajuste:}
\[
|p_{\text{dev}}| \geq \bar{p}(z) = \sqrt{\frac{2z}{\eta-1}}
\]

\textbf{Curva de Lucas:}
\[
y = b\,(p - \E[p])
\]

\textbf{Curva de Phillips NK (Calvo):}
\[
\pi_t = \beta\,\E_t[\pi_{t+1}] + \kappa\,y_t,\quad
\kappa \equiv \frac{\alpha[1-(1-\alpha)\beta]\phi}{1-\alpha}
\]
onde $\phi = \theta + \gamma - 1$.

\textbf{CPNK com inércia (Christiano et al.):}
\[
\pi_t = \gamma\pi_{t-1} + (1-\gamma)\E_t[\pi_{t+1}] + \chi y_t
\]

\columnbreak

\textbf{\textcolor{cyanrbc}{BLOCO MODELO NK CANÔNICO}}

\medskip
\textbf{IS dinâmica:}
\[
y_t = \E_t[y_{t+1}] - \frac{1}{\theta}\,r_t + u_t^{IS}
\]

\textbf{CPNK (3-equação):}
\[
\pi_t = \beta\,\E_t[\pi_{t+1}] + \kappa\,y_t + u_t^{\pi}
\]

\textbf{Regra de Taylor (em juros reais):}
\[
r_t = \phi_\pi\,\E_t[\pi_{t+1}] + \phi_y\,\E_t[y_{t+1}] + u_t^{MP}
\]

\textbf{Princípio de Taylor (determinação):}
\[
\phi_\pi > 1 \;\;\text{(Blanchard-Kahn)}
\]

\textbf{Modelo de Fischer ($\phi$ = inv.\ rigidez real):}
\[
p_t = \E_{t-2}m_t + \frac{\phi}{1+\phi}[\E_{t-1}m_t - \E_{t-2}m_t]
\]
\[
y_t = \frac{1}{1+\phi}[\E_{t-1}m_t - \E_{t-2}m_t] + [m_t - \E_{t-1}m_t]
\]

\textbf{Modelo de Taylor:}
\[
y_t = \lambda\,y_{t-1} + \frac{1+\lambda}{2}\,u_t
\]

\textbf{Calibração NK padrão:}\\[2pt]
$\beta=0{,}99$, $\kappa=0{,}1$, $\phi_\pi=1{,}5$, $\phi_y=0{,}5$,\\
$\theta=1{,}0$, $\alpha=0{,}75$ (Calvo)

\end{multicols}

\rule{\linewidth}{0.6pt}

\vspace{0.5em}
\begin{provabox}[title={Fórmula mais cobrada: inclinação da CPNK de Calvo}]
A inclinação $\kappa$ mede o pass-through do hiato do produto para a inflação.
Quanto maior $\alpha$ (fração de firmas \emph{que} ajustam), maior $\kappa$ e mais íngreme a CPNK.
Quanto maior a rigidez real ($\phi$ menor), menor $\kappa$ --- a CPNK fica mais plana e a
política monetária precisa de maiores variações de juros para controlar a inflação.
\end{provabox}

% --------------------------------------------------------------------------
% Sumário
% --------------------------------------------------------------------------
\newpage
\tableofcontents
\newpage

% ==========================================================================
%  PARTE I — RIGIDEZ NOMINAL E COMPETIÇÃO MONOPOLÍSTICA (Q1-Q4)
% ==========================================================================

\section*{\large\color{cyanrbc} Parte I --- Rigidez Nominal e Preços (Q1--Q4)}
\addcontentsline{toc}{section}{Parte I --- Rigidez Nominal e Preços (Q1--Q4)}

Os modelos de rigidez nominal surgidos a partir da década de 1970 (Mankiw, 1985;
Ball \& Romer, 1990; Calvo, 1983) mostram como pequenas fricções na formação de
preços em nível microeconômico geram não-neutralidade da moeda no curto prazo.
O elemento central é a competição monopolística à la Dixit-Stiglitz, que fornece
microfundamentos para o markup e para o produto de equilíbrio.

\bigskip
\setcounter{section}{0}

% ==========================================================================
\section{Rigidez nominal e efeitos da política monetária}

\subsection*{Enunciado}
Com base em modelos surgidos a partir da década de 70, a ocorrência de alguma forma
de rigidez nominal na economia é essencial para: (a) efeitos nominais da PM; (b) efeitos
reais de longo prazo de choques nos dispêndios; (c) coordenação entre PM e fiscal;
(d) otimização dos agentes; \textbf{(e) efeitos reais de curto prazo dos choques monetários}.

\begin{resultbox}[title={Resposta: alternativa (e)}]
A rigidez nominal é a condição necessária para que choques monetários tenham efeitos
\textbf{reais no curto prazo}. Sem ela, os agentes ajustariam preços imediatamente e a
moeda seria neutra em qualquer horizonte. As demais alternativas não requerem rigidez
nominal: efeitos nominais da PM existem mesmo com preços flexíveis; efeitos reais de
longo prazo dependem de fatores reais (capital, tecnologia); coordenação PM-fiscal é
questão de regra fiscal; otimização dos agentes é compatível com preços flexíveis.
\end{resultbox}

% ==========================================================================
\section{Otimização da família: consumo e oferta de trabalho ótimos}

\subsection*{Enunciado}
Seja $U = C^{1-\theta} - \tfrac{1}{\gamma}L^\gamma$, $\gamma>1$, $\theta>0$, com restrição
$C = WL$. Calcule $C^*$ e $L^*$.

\subsection*{Resolução}

Formamos o lagrangiano com multiplicador $\lambda$:
\begin{equation}
\mathcal{L} = C^{1-\theta} - \frac{1}{\gamma}L^\gamma + \lambda(WL - C).
\label{eq:lagr2}
\end{equation}

\textbf{CPO em $C$:}
\[
\frac{\partial\mathcal{L}}{\partial C} = (1-\theta)C^{-\theta} - \lambda = 0
\;\Longrightarrow\;
\lambda = (1-\theta)C^{-\theta}. \tag{2}
\]

\textbf{CPO em $L$:}
\[
\frac{\partial\mathcal{L}}{\partial L} = -L^{\gamma-1} + \lambda W = 0
\;\Longrightarrow\;
\lambda = \frac{L^{\gamma-1}}{W}. \tag{3}
\]

\textbf{Igualando (2) e (3):}
\begin{equation}
(1-\theta)C^{-\theta} = \frac{L^{\gamma-1}}{W}.
\label{eq:foc_equal}
\end{equation}

\begin{resultbox}[title={Consumo e oferta de trabalho ótimos}]
\textbf{Resolvendo \eqref{eq:foc_equal} para $C$:}
\[
\boxed{C^* = \left[\frac{W(1-\theta)}{L^{\gamma-1}}\right]^{1/\theta}}
\]

\textbf{Resolvendo \eqref{eq:foc_equal} para $L$:}
\[
\boxed{L^* = \left[W(1-\theta)C^{-\theta}\right]^{1/(\gamma-1)}}
\]
\end{resultbox}

\begin{provabox}[title={Intuição econômica}]
O consumo ótimo $C^*$ aumenta com o salário $W$ (maior renda) e diminui com $\gamma$
(maior desutilidade do trabalho reduz o equilíbrio de trabalho e, portanto, a renda).
A oferta ótima $L^*$ aumenta com $W$ (efeito substituição) e com $(1-\theta)$
(menor aversão ao risco aumenta a utilidade marginal do consumo, incentivando mais trabalho).
\end{provabox}

% ==========================================================================
\section{Produto de equilíbrio com competição monopolística}

\subsection*{Enunciado}
Seja $Y = \left[\tfrac{\eta-1}{\eta}\right]^{1/(\gamma-1)}$. Interprete os efeitos de
aumentos em $\eta$ e $\gamma$ sobre $Y^*$.

\subsection*{Derivação e interpretação}

O produto de equilíbrio resulta da igualação entre a demanda do consumidor (derivada
do problema de maximização de utilidade) e a regra de markup da firma monopolista:
\[
P_i^* = \frac{\eta}{\eta-1}\,MC_i = \mu\,MC_i, \qquad \mu \equiv \frac{\eta}{\eta-1}.
\]
No equilíbrio simétrico ($P_i = P$), combinando as CPOs da família e da firma,
obtém-se:
\[
Y^* = \left[\frac{\eta-1}{\eta}\right]^{1/(\gamma-1)}.
\]

\begin{resultbox}[title={Efeitos de $\eta$ e $\gamma$}]
\textbf{Aumento de $\eta$:} reduz o markup $\mu = \eta/(\eta-1)$, diminuindo o poder
de mercado das firmas. Isso eleva $Y^*$ --- mais concorrência $\Rightarrow$ menor
distorção $\Rightarrow$ maior produto de equilíbrio.

\[
\eta \uparrow \;\Rightarrow\; \mu \downarrow \;\Rightarrow\; \frac{\eta-1}{\eta} \uparrow
\;\Rightarrow\; Y^* \uparrow
\]

\textbf{Aumento de $\gamma$:} reduz o expoente $1/(\gamma-1)$, diminuindo a
elasticidade da oferta de trabalho a variações salariais. Isso contrai $Y^*$.

\[
\gamma \uparrow \;\Rightarrow\; \frac{1}{\gamma-1} \downarrow \;\Rightarrow\; Y^* \downarrow
\]
\end{resultbox}

% ==========================================================================
\section{Nível ótimo de preços e efeito de $\gamma$}

\subsection*{Enunciado}
Dada $Y = M/P$ (demanda agregada simples), deduza o nível de preços de equilíbrio.
Coeteris paribus, o que ocorre com $P$ após elevação de $\gamma$? Interprete.

\subsection*{Resolução}

Substituindo $Y^*$ em $P = M/Y^*$:

\begin{resultbox}[title={Nível de preços de equilíbrio}]
\[
\boxed{P^* = \frac{M}{\left[\dfrac{\eta-1}{\eta}\right]^{1/(\gamma-1)}}}
\]
\end{resultbox}

\textbf{Efeito de $\gamma \uparrow$:} Da questão anterior, $\gamma \uparrow \Rightarrow Y^* \downarrow$.
Com $M$ fixo (ceteris paribus), o excesso de demanda nominal eleva o nível de preços:
\[
\gamma \uparrow \;\Rightarrow\; Y^* \downarrow \;\Rightarrow\; P^* = \frac{M}{Y^*} \uparrow.
\]

\begin{provabox}[title={Interpretação estrutural}]
A elevação de $\gamma$ representa maior desutilidade marginal do trabalho, reduzindo
a oferta de trabalho de equilíbrio e, portanto, o produto. Dada a oferta de moeda
pelo Banco Central, o ajuste recai inteiramente sobre os preços: a economia precisa
de um nível de preços mais alto para equilibrar o mercado de bens. Esse mecanismo
é análogo a um choque negativo de oferta agregada na curva AS.
\end{provabox}

% ==========================================================================
%  PARTE II — INFORMAÇÃO INCOMPLETA E CRÍTICA DE LUCAS (Q5-Q7)
% ==========================================================================

\newpage
\section*{\large\color{cyanrbc} Parte II --- Informação Incompleta e Crítica de Lucas (Q5--Q7)}
\addcontentsline{toc}{section}{Parte II --- Informação Incompleta e Crítica de Lucas (Q5--Q7)}

O modelo de informação incompleta de Lucas (1973) mostra que mesmo agentes
otimizadores e racionais podem ser enganados por choques monetários, pois não
conseguem distinguir, em tempo real, variações de preços relativos de variações
do nível geral de preços.

\bigskip

% ==========================================================================
\section{Modelo de informação incompleta de Lucas}

\subsection*{Enunciado}
A firma decide $y_i = [1/(\gamma-1)]\,\E[r_i \mid p_i]$, onde $r_i = p_i - p$.
Explique como flutuações de produto no curto prazo resultam de falta de informação
e como a rigidez de preços pode ser deduzida desse modelo.

\begin{resultbox}[title={Mecanismo de extração de sinal}]
A firma \textbf{não observa} $p$ (nível geral) no momento da decisão de produção;
ela apenas observa seu próprio preço $p_i$. Diante de um choque positivo de demanda:
\begin{enumerate}[leftmargin=1.5em, itemsep=2pt]
  \item A demanda pelo produto $i$ sobe $\Rightarrow$ $p_i \uparrow$.
  \item A firma não sabe se $p_i \uparrow$ reflete $r_i = p_i - p \uparrow$
    (aumento \emph{relativo}) ou apenas $p \uparrow$ (inflação geral).
  \item Ela atribui parte do aumento a $r_i \uparrow$ $\Rightarrow$ eleva $y_i$.
\end{enumerate}
\end{resultbox}

\textbf{Rigidez de preços endógena:} como a firma não ajusta seu preço de forma
instantânea e completa ao nível geral --- pois não dispõe dessa informação ---
os preços individuais respondem menos do que o necessário ao choque monetário,
gerando a aparência de rigidez nominal mesmo sem custos de menu explícitos.

% ==========================================================================
\section{Curva de Lucas: choque monetário e convergência de longo prazo}

\subsection*{Enunciado}
Seja $y = b(p - \E[p])$, $b>0$. Interprete a dinâmica do produto após um choque
monetário positivo. Explique a convergência de longo prazo.

\subsection*{Mecanismo de curto e longo prazo}

\textbf{Curto prazo:} Um choque $m > \E[m]$ eleva $p$ acima de $\E[p]$:
\[
m \uparrow \;\Rightarrow\; p > \E[p] \;\Rightarrow\; y = b(p - \E[p]) > 0.
\]
As firmas interpretam parte da alta de $p_i$ como aumento de preço relativo,
expandindo a produção.

\begin{resultbox}[title={Convergência de longo prazo}]
No longo prazo, via expectativas racionais e aprendizado, os agentes atualizam
$\E[m]$ e $\E[p]$ para os novos níveis:
\[
\E[m] \uparrow \;\text{ e }\; \E[p] \uparrow \;\Rightarrow\; p = \E[p] \;\Rightarrow\; y = 0.
\]
Uma vez que a nova regra de política monetária é assimilada, a moeda é
\textbf{neutra no longo prazo}: o produto retorna ao nível normal (desvio nulo).
A correlação positiva entre $p$ e $y$ desaparece --- consistente com a
dicotomia clássica de longo prazo.
\end{resultbox}

% ==========================================================================
\section{Crítica de Lucas}

\subsection*{Enunciado}
Um BC estuda mudar a regra de PM com base em séries temporais. Interprete a
eficácia da medida a partir da Crítica de Lucas.

\begin{resultbox}[title={Crítica de Lucas (1976)}]
As relações econométricas estimadas com dados passados são válidas apenas
\textbf{condicional ao estado de expectativas} prevalecente naquele período.
Uma mudança de \emph{regra} de PM altera estruturalmente as expectativas dos agentes,
invalidando as previsões feitas com o modelo antigo.
\end{resultbox}

\begin{provabox}[title={Implicação para política econômica}]
Para que as previsões baseadas em séries históricas sejam válidas, seria necessário
que a mudança de PM \textbf{não fosse percebida como estrutural} pelos agentes.
Mas se os agentes são racionais (Muth, 1961), eles antecipam a nova regra e ajustam
suas expectativas imediatamente --- tornando o modelo histórico obsoleto.
Isso é a base teórica para a independência do BC e para a comunicação de políticas
críveis via ancoragem de expectativas.
\end{provabox}

% ==========================================================================
%  PARTE III — MODELOS DE RIGIDEZ: FISCHER, TAYLOR, CALVO (Q8-Q10)
% ==========================================================================

\newpage
\section*{\large\color{cyanrbc} Parte III --- Modelos de Rigidez de Preços (Q8--Q10)}
\addcontentsline{toc}{section}{Parte III --- Modelos de Rigidez (Q8--Q10)}

Os modelos de Fischer (1977), Taylor (1980) e Calvo (1983) formalizaram a rigidez
nominal com microfundamentos distintos. Em todos eles, o parâmetro chave é o
\textbf{grau de rigidez real} ($\phi$ ou $\alpha$), que determina o pass-through
de choques monetários para preços versus produto.

\bigskip

% ==========================================================================
\section{Modelos de Fischer e Taylor: rigidez e efeitos reais da PM}

\subsection*{Enunciado}
Comente sobre o grau de rigidez de preços (a partir de $\phi$) e os efeitos reais
da PM nos modelos de: (a) Fischer; (b) Taylor.

\subsection*{(a) Modelo de Fischer}

O modelo assume contratos de dois períodos: os preços do período $t$ são fixados em
$t-2$ (firmas do grupo 1) e em $t-1$ (firmas do grupo 2). As equações de equilíbrio são:
\begin{align}
p_t &= \E_{t-2}m_t + \frac{\phi}{1+\phi}\left[\E_{t-1}m_t - \E_{t-2}m_t\right], \\
y_t &= \frac{1}{1+\phi}\left[\E_{t-1}m_t - \E_{t-2}m_t\right] + \left[m_t - \E_{t-1}m_t\right].
\end{align}

\begin{resultbox}[title={Papel de $\phi$ no modelo de Fischer}]
$\phi$ é o \textbf{inverso do grau de rigidez real}: maior $\phi$ = menor rigidez.
\begin{itemize}[leftmargin=1.5em, itemsep=2pt]
  \item $\phi \uparrow$: os preços respondem mais às revisões de expectativas
    $[\E_{t-1}m_t - \E_{t-2}m_t]$, pois as firmas têm maior flexibilidade real.
  \item $\phi \uparrow$: o coeficiente do produto $1/(1+\phi)$ \emph{diminui} --- menor
    efeito real da PM sobre o produto.
  \item Intuição: quando os preços são mais flexíveis, eles absorvem mais os choques,
    deixando menos espaço para efeitos reais.
\end{itemize}
\end{resultbox}

\subsection*{(b) Modelo de Taylor}

O modelo de Taylor assume contratos escalonados de dois períodos. O resultado
macroeconômico é:
\[
y_t = \lambda\,y_{t-1} + \frac{1+\lambda}{2}\,u_t,
\]
onde $\lambda$ mede a \textbf{sensibilidade dos preços correntes aos preços passados}
(rigidez real e persistência do produto) e $u_t$ é o choque de demanda.

\begin{resultbox}[title={Papel de $\lambda$ e $\phi$ no modelo de Taylor}]
\begin{itemize}[leftmargin=1.5em, itemsep=2pt]
  \item $\lambda$ e $\phi$ têm relação inversa: $\lambda \uparrow \Leftrightarrow \phi \downarrow$
    (mais rigidez real $\Rightarrow$ maior persistência do produto).
  \item $\phi \uparrow$ (menos rigidez): menor $\lambda$, menor persistência de $y_t$
    e menor efeito de $u_t$ sobre o produto.
  \item O modelo de Taylor gera \textbf{persistência endógena} dos ciclos de negócios
    via escalonamento de contratos, mesmo sem choques persistentes.
\end{itemize}
\end{resultbox}

% ==========================================================================
\section{Modelo de Calvo: papel de $\theta$ e $\alpha$ na CPNK}

\subsection*{Enunciado}
Na CPNK de Calvo $\pi_t = \kappa y_t + \beta\E_t\pi_{t+1}$, com
$\kappa \equiv \alpha[1-(1-\alpha)\beta]\phi/(1-\alpha)$ e $\phi = \theta + \gamma - 1$,
comente o papel de: (a) $\theta$ (aversão ao risco); (b) $\alpha$ (fração de firmas
que ajustam).

\begin{resultbox}[title={(a) Papel de $\theta$ — aversão ao risco}]
$\theta$ entra no slope via $\phi = \theta + \gamma - 1$. Um aumento de $\theta$:
\[
\theta \uparrow \;\Rightarrow\; \phi \uparrow \;\Rightarrow\; \kappa \uparrow.
\]
Maior aversão ao risco \textbf{reduz a rigidez real} (firmas querem ajustar preços
mais rapidamente para evitar perdas), tornando a CPNK \textbf{mais íngreme}.
Assim, variações do hiato do produto têm impacto inflacionário \textbf{maior} quando
$\theta$ é alto --- a PM precisa ser mais agressiva para controlar a inflação.
\end{resultbox}

\begin{resultbox}[title={(b) Papel de $\alpha$ — fração de firmas que ajustam}]
$\alpha$ aparece tanto no numerador quanto no denominador de $\kappa$:
\[
\alpha \uparrow \;\Rightarrow\; \kappa \uparrow \quad \text{(CPNK mais íngreme)}.
\]
Intuitivamente: se mais firmas podem ajustar preços em cada período, a inflação
responde mais rapidamente a variações na atividade econômica.
O caso limite $\alpha \to 1$ (todas as firmas ajustam) recupera a curva de oferta
agregada com preços totalmente flexíveis.
\end{resultbox}

% ==========================================================================
\section{Calvo (1983) vs Christiano et al. (2005): inércia inflacionária}

\subsection*{Enunciado}
Explique a diferença entre os modelos e como Christiano et al. (2005) contorna o
problema da desinflação sem recessão do Calvo original.

\subsection*{Limitação do modelo de Calvo puro}

No modelo de Calvo puro, a CPNK é estritamente forward-looking:
$\pi_t = \beta\E_t\pi_{t+1} + \kappa y_t$.
Um processo de desinflação crível ($\E_t\pi_{t+1} \downarrow$) leva a
$\pi_t \downarrow$ \textbf{sem} necessidade de recessão ($y_t$ não precisa cair).
Isso contradiz a evidência empírica de que desinflações geram \textbf{custos de sacrifício}.

\subsection*{Solução de Christiano et al. (2005): indexação}

O modelo assume contratos com indexação à inflação passada. A CPNK resultante é:
\begin{equation}
\pi_t = \gamma\pi_{t-1} + (1-\gamma)\E_t[\pi_{t+1}] + \chi y_t, \quad 0 \leq \gamma \leq 1.
\end{equation}

\begin{resultbox}[title={Mecanismo de geração de custos de desinflação}]
Quando $\gamma > (1-\gamma)$ (o passado pesa mais que as expectativas futuras), uma
desinflação $\pi_{t-1} > \pi_t > \E_t\pi_{t+1}$ exige $y_t < 0$.
\end{resultbox}

\begin{provabox}[title={Exemplo numérico ($\gamma=0{,}8$, $\chi=0{,}9$)}]
Dados: $\E_t\pi_{t+1} = 1{,}5\%$, $\pi_{t-1} = 6\%$, $\pi_t = 3{,}5\%$.
\[
3{,}5 = 0{,}8 \times 6 + 0{,}2 \times 1{,}5 + 0{,}9\,y_t
\;\Rightarrow\;
y_t = \frac{3{,}5 - 4{,}8 - 0{,}3}{0{,}9} = \frac{-1{,}6}{0{,}9} \approx -1{,}7\%.
\]
A desinflação custou aproximadamente $1{,}7\%$ de hiato negativo do produto ---
consistente com evidências empíricas de razões de sacrifício positivas.
\end{provabox}

% ==========================================================================
%  PARTE IV — MODELO NK CANÔNICO (Q11)
% ==========================================================================

\newpage
\section*{\large\color{cyanrbc} Parte IV --- Modelo NK Canônico (Q11)}
\addcontentsline{toc}{section}{Parte IV --- Modelo NK Canônico (Q11)}

O modelo Novo-Keynesiano canônico de três equações (Woodford, 2003; Galí, 2008)
oferece uma representação compacta da demanda agregada (IS), oferta agregada (CPNK)
e comportamento do BC (Regra de Taylor). Ele é o principal framework para análise
de política monetária em economias com rigidez nominal.

\bigskip

% ==========================================================================
\section{Modelo NK canônico: aversão ao risco, choques de oferta e PM}

\subsection*{Enunciado}
O modelo canônico NK:
\begin{align}
y_t &= \E_t[y_{t+1}] - \frac{1}{\theta}\,r_t + u_t^{IS},\quad \theta > 0, \\
\pi_t &= \beta\,\E_t[\pi_{t+1}] + \kappa\,y_t + u_t^{\pi},\quad 0<\beta<1,\;\kappa>0, \\
r_t &= \phi_\pi\,\E_t[\pi_{t+1}] + \phi_y\,\E_t[y_{t+1}] + u_t^{MP},\quad \phi_\pi>0,\;\phi_y\geq 0.
\end{align}
Comente: (i) como $\theta$ afeta as relações estruturais; (ii) como choques negativos
de oferta influenciam a condução de PM.

\subsection*{(i) Papel de $\theta$ (aversão ao risco)}

\textbf{Efeito na IS:} $\theta$ é a inversa da elasticidade intertemporal de substituição
(EIS). Maior $\theta$ $\Rightarrow$ famílias respondem menos a variações de juros reais:
\[
\theta \uparrow \;\Rightarrow\; \left|\frac{\partial y_t}{\partial r_t}\right| = \frac{1}{\theta} \downarrow.
\]

\begin{resultbox}[title={Efeitos de $\theta$ no sistema NK}]
\begin{enumerate}[leftmargin=1.5em, itemsep=3pt]
  \item \textbf{IS mais inelástica:} a PM precisa de variações de juros \emph{maiores}
    para eliminar o hiato do produto. Porém, o custo social é nulo (só move $r_t$).
  \item \textbf{CPNK mais íngreme:} $\theta \uparrow \Rightarrow \phi = \theta+\gamma-1 \uparrow
    \Rightarrow \kappa \uparrow$. Choques de demanda têm maior impacto inflacionário.
  \item \textbf{Regra de Taylor mais agressiva:} $\phi_\pi$ e $\phi_y$ ótimos aumentam
    com $\theta$ e $\kappa$ (e diminuem com $\beta$), pois o BC precisa reagir mais
    fortemente para compensar a maior rigidez e o maior pass-through.
\end{enumerate}
\end{resultbox}

\subsection*{(ii) Choques negativos de oferta e condução de PM}

Um choque negativo de oferta corresponde a $u_t^\pi > 0$ na CPNK. Isso cria o
\textbf{dilema fundamental} da PM sob rigidez nominal:

\begin{resultbox}[title={Trade-off estagflacionário}]
\[
u_t^\pi > 0 \;\Rightarrow\; \begin{cases}
\pi_t \uparrow & (\text{pressão inflacionária}) \\
y_t \downarrow & (\text{custo real se PM reagir})
\end{cases}
\]
Para zerar o desvio de inflação, o BC eleva $r_t$ $\Rightarrow$ $y_t < 0$
(custo social positivo = recessão). Não é possível estabilizar $\pi_t$ e $y_t$
simultaneamente após um choque de custo-push.
\end{resultbox}

\begin{provabox}[title={Estratégias do BC diante de choques de oferta}]
\textbf{Opção 1 --- Desinflação imediata:} Elevar $r_t$ para $\pi_t = 0$.
Custo: hiato negativo $y_t < 0$ (recessão). Benefício: credibilidade preservada.

\textbf{Opção 2 --- Acomodação transitória:} Tolerar $\pi_t > 0$ temporariamente
(bandas de tolerância), esperando a reversão natural do choque. Justificada quando
$\rho_\pi$ (persistência) é baixa.

\textbf{Opção 3 --- Desinflação gradual:} Elevar $r_t$ progressivamente ao longo
de vários períodos. Minimiza o custo de recessão em $t$, mas prolonga o desvio
inflacionário.

Quando $\rho_\pi$ é muito alto, o BC tende a preferir a Opção 1 ou 3 para evitar
a perda de controle da dinâmica inflacionária e o risco de \emph{desancoragem}
das expectativas.
\end{provabox}

\vspace{1em}
\begin{formulabox}[title={Conexão com o modelo computacional (ch07\_nk.py)}]
O modelo NK canônico das três questões acima está completamente implementado em
\texttt{ch07\_dsge\_nk/ch07\_nk.py} (classe \texttt{NKModel}).
A solução analítica por coeficientes indeterminados ($\psi_{xv}$, $\psi_{\pi v}$,
$\psi_{xu}$, $\psi_{\pi u}$) reproduz exatamente os efeitos qualitativos discutidos.
As IRFs para choques de demanda ($v_t$) e oferta ($u_t$) confirmam:
$\psi_{xv} < 0$, $\psi_{\pi v} < 0$ (aperto monetário contrai $y$ e $\pi$);
$\psi_{xu} < 0$, $\psi_{\pi u} > 0$ (choque de custo-push gera estagflação).
\end{formulabox}

\end{document}
